\fichatitulo{17 · CORPUS}{Texto para Discussão 355 · 2025}{Plenário, palanque e estúdio}
{\small Blanco. \textit{Plen{\'a}rio, Palanque, Est{\'u}dio: Discursos no Plen{\'a}rio do Senado Federal entre 2007 e 2024}. Texto para Discussão 355 · 2025. \href{https://www12.senado.leg.br/publicacoes/estudos-legislativos/tipos-de-estudos/textos-para-discussao/td355}{Fonte consultada}.\par}

\campo{Problema e objetivo}
Como se transformaram os discursos no Senado entre 2007 e 2024?

\campo{Principal contribuição · síntese interpretativa}
Relaciona evolução dos pronunciamentos, dinâmica do plenário e comunicação audiovisual por meio de indicadores descritivos.

\campo{Método / natureza da evidência}
Análise de discursos e metadados, com processamento computacional e exploração temática.

\campo{Resultado ou argumento principal}
Identifica redução de frequência, duração e interação, com recuperação parcial posterior; discute mudanças na função comunicativa dos pronunciamentos.

\campo{Excertos-chave · seleção literal}
\begin{itemize}
\excerto{1}{As medidas aqui apresentadas são essencialmente descritivas}{Discussão.}
\excerto{2}{correlação não implica causação}{Discussão.}
\end{itemize}

\campo{Limitações e análise crítica}
Autor: as medidas são descritivas e correlação não implica causação. Análise da ficha: explicações envolvendo redes sociais exigem desenhos próprios; o corpus não deve ser presumido idêntico ao da dissertação.

\campo{Aproveitamento na dissertação · proposta}
Contextualização do corpus: documentar temporalidade, seleção dos discursos e mudanças de formato.

\registro{Fonte consultada nesta preparação: Resumo, apresentação metodológica, discussão e conclusão. Chave: \texttt{blanco2025plenarioPalanqueEstudio}.}
